
%%%%%%%%%%%%%%%%%%%%%%%%%%%%%%%%%%%%%%%%%%%%%%%%%%%%%%%%%%%%%
\section{模型的评价}

\subsection{模型的优点}
\begin{itemize}[itemindent=2em]
\item \textbf{机理为主、外推可靠：}以Deutsch方程等物理机理模型为主框架，从历史数据只拟合待定系数，外推时靠物理单调性约束保证趋势合理，避免了纯数据驱动在$C_{out}$方差极小时外推不可靠的问题。
\item \textbf{工况自适应寻优：}K-Means工况划分使不同入口条件下的最优参数独立求解，避免了"一刀切"统一参数的弊端，高浓度与低浓度工况策略差异显著，符合工程实际。
\item \textbf{排放阈值参数化设计：}问题四复用问题二的优化框架，仅需传入不同$C_{limit}$即可求解，代码与模型结构一致性好，便于扩展到其他排放标准。
\item \textbf{多起点+全局备选保证收敛：}SLSQP多起点启动+差分进化备选，有效应对非凸可行域，6个工况均成功收敛且约束满足。
\end{itemize}